% Per-Layer Delta Features and Manifold Signatures in
%   KV-Cache Confabulation Detection
%
% Companion to: Spectral Shape Features paper (shape beats threshold)
% This paper: layer averaging destroys signal, delta beats endpoint,
%   trajectory geometry shows manifold structure

\documentclass[11pt,a4paper]{article}
\usepackage{amsmath,amssymb,booktabs,graphicx,hyperref}
\usepackage{natbib}
\usepackage{multirow}

\title{Per-Layer Delta Features and Manifold Signatures\\
in KV-Cache Confabulation Detection}

\author{
Lyra\textsuperscript{1*}\quad
CC\textsuperscript{1}\quad
Thomas Edrington\textsuperscript{1}\quad
Dwayne Wilkes\textsuperscript{1,2}\\[4pt]
\textsuperscript{1}Liberation Labs\quad
\textsuperscript{2}Sentient Futures\\
\textsuperscript{*}Lead AI author. Corresponding human author: T.~Edrington.
}

\date{Draft --- \today}

\begin{document}
\maketitle

% ===================================================================
\begin{abstract}
% ===================================================================

Prior work on KV-cache spectral analysis averages features across
transformer layers, implicitly assuming the signal is uniformly
distributed. We show this assumption is costly. Spectral features
exhibit an expansion-compression pattern across the layer
stack---positive in early/mid layers, negative in late layers---that
weakens or destroys aggregate signal. On Qwen2.5-7B, per-layer
confabulation detection peaks at $d = 2.46$ (layer~23), 37\% stronger
than the all-layer average ($d = 1.80$). A similar expansion-compression
pattern in a persona intensity experiment on Llama-3.1-8B produces
$d = -5.15$ at layer~29 (spectral entropy; note: 1.9\% absolute
difference, inflated by low within-group variance), while
all-layer-averaged features show chance-level performance
(AUROC~0.520). Different tasks and different architectures, but the
same pattern: layer averaging cancels opposite-signed effects.

We introduce delta features (generation minus encoding phase), which
achieve $d = 2.35$ ($p < 0.001$) for confabulation detection on
Qwen2.5-7B. The delta uses information from both phases and is not
directly comparable to single-phase endpoints ($d = 1.80$); the
improvement may partly reflect access to more information rather than
a purely superior signal. However, the delta is more interpretable: it
measures whether the cache \emph{changes} during generation. Grounded
responses expand the cache (delta stable\_rank $+0.44$);
confabulation leaves it flat ($+0.03$).

Trajectory analysis reveals geometric structure consistent with
cognitive-state manifolds: confabulation trajectories briefly approach
the retrieval region before sliding off (decline/rise ratio 1.91 vs
0.40--0.56 for grounded conditions).

\end{abstract}

% ===================================================================
\section{Introduction}
\label{sec:intro}
% ===================================================================

Companion work~\citep{lyra_spectral} established that
threshold-independent spectral shape features (stable rank, singular
value kurtosis, participation ratio) outperform Marchenko-Pastur
features for confabulation detection from the KV cache, achieving
AUROC~0.767 on Qwen2.5-7B-Instruct. That analysis aggregated
features by computing the mean across all transformer layers---a
common default in cache-based analysis.

This paper shows that default is wrong. The confabulation signal is
not uniform across layers. It exhibits an \emph{expansion-compression
  cycle}: spectral features increase (expand) in early-to-mid layers
and decrease (compress) in late layers. When averaged, these
opposite-signed effects cancel, destroying signal that per-layer
analysis recovers.

The consequences can be severe. On Llama-3.1-8B-Instruct, all-layer
averaging of confabulation detection features produces
AUROC~0.520---indistinguishable from chance. Yet in a separate persona
intensity experiment on the same architecture, a single late layer
(L29) yields $d = -5.15$ for spectral entropy (though the absolute
difference is small---1.9\%---and $d$ is inflated by low
within-group variance). This suggests strong layer-specific signal
that averaging renders invisible.

We make three contributions:

\begin{enumerate}
\item \textbf{Per-layer analysis recovers signal that averaging
    weakens or destroys.} A consistent expansion-compression pattern
  appears across experiments on Qwen (28 layers) and Llama (32
  layers), with crossover at approximately 50\% depth. Different
  tasks and features were measured on each architecture
  (Section~\ref{sec:limitations}), so we report this as a consistent
  pattern rather than a confirmed universal property.

\item \textbf{Delta features provide a strong, interpretable
    detector.} The difference between generation-phase and
  encoding-phase stable rank ($d = 2.35$, $p < 0.001$) exceeds any
  single-phase measurement ($d = 1.80$), though the comparison is not
  strictly fair because the delta uses information from both phases.
  Its interpretive advantage is clear: it directly measures whether
  the cache changed during generation.

\item \textbf{Trajectory geometry provides evidence for manifold
    structure.} Token-by-token trajectories through feature space show
  that confabulation follows a distinctive ``approach and slide''
  pattern (rise-fall ratio 1.91) absent from grounded conditions
  (0.40). Pairwise condition distances exhibit non-linear structure
  (detour ratio 1.24), consistent with curved cognitive-state
  manifolds.
\end{enumerate}

% ===================================================================
\section{Related Work}
\label{sec:related}
% ===================================================================

\paragraph{KV-cache spectral analysis.}
Prior work in this program established Marchenko-Pastur corrected SVD
for confabulation detection~\citep{lyra_p2}, emotion-vector
steering~\citep{lyra_p6}, cache obfuscation
defense~\citep{lyra_p7}, and user-state decoding from encoding-phase
cache~\citep{lyra_p4}. The companion spectral shape
paper~\citep{lyra_spectral} introduced threshold-free features and the
retrieval-engagement interpretation. All prior work used all-layer
averaging.

\paragraph{Layer-specific representations.}
It is well established that different transformer layers encode
different levels of abstraction~\citep{tenney2019bert}. Early layers
capture syntactic structure; later layers capture
semantics~\citep{jawahar2019does}. Our finding that spectral features
have opposite signs at different depths is consistent with this
hierarchy but has not been previously reported for KV-cache geometry.

\paragraph{Manifold representations.}
Recent work demonstrates that neural network representations live on
curved low-dimensional manifolds~\citep{goodfire2026}. Linear
steering methods fail by pushing representations into off-manifold
``voids.'' Sparse autoencoder features tend to ``shatter'' manifolds
into local fragments. Our spectral shape features, which measure
global distributional properties, may capture manifold geometry more
directly.

% ===================================================================
\section{Methods}
\label{sec:methods}
% ===================================================================

\subsection{Data Sources}

We analyze two datasets from the KV-cache research program:

\begin{enumerate}
\item \textbf{Honesty signal experiment} ($n = 130$ trials,
  Qwen2.5-7B-Instruct): five conditions (honest-common, honest-rare,
  confab, deceptive-user, boundary), fixed 100-token generation,
  per-layer features at all 28 layers, trajectory checkpoints every 10
  tokens. Source: \citet{lyra_spectral}.

\item \textbf{Persona intensity experiment} ($n = 50$ prompts $\times$
  5 levels $\times$ 3 models): Qwen2.5-7B (28 layers),
  Llama-3.1-8B (32 layers), Mistral-7B (32 layers). Five persona
  intensity levels (L0 through L3, plus L0.5 style-only control).
  Per-layer features at all layers. Source: CC's cross-architecture
  analysis, red-teamed.
\end{enumerate}

\subsection{Feature Extraction}

All features extracted using \texttt{lyra\_features.py}, computing
per-layer singular values via GPU SVD
(\texttt{torch.linalg.svdvals}). Features include stable rank,
participation ratio, sv\_kurtosis, spectral entropy, condition number,
and MP/GD threshold-based features.

\paragraph{Per-layer features.} Raw feature values at each layer,
without aggregation. This is the key methodological change from prior
work.

\paragraph{Delta features.} For each trial, delta = generation-phase
feature $-$ encoding-phase feature (all-layer mean for both phases).
This measures how much the cache \emph{changes} during generation.

\paragraph{Trajectory features.} Stable rank extracted at 10
equispaced checkpoints during generation (every 10 tokens for
100-token generation). The trajectory is a curve through feature
space.

\subsection{Analysis}

\begin{itemize}
\item Per-layer effect sizes: Cohen's $d$ between conditions at each
  layer, with paired $t$-test $p$-values.
\item Delta comparison: two-sample $t$-test on delta stable\_rank
  between conditions.
\item Trajectory shape classification: each trajectory classified as
  ``rise and stay'' (peak in final 20\% of generation) or ``rise and
  fall'' (peak before 80\%, decline $>$ 30\% of rise). Decline/rise
  ratio computed for rise-and-fall trajectories.
\item Distance matrix: z-scored Euclidean distance between condition
  centroids on 6-feature endpoint vectors.
\item Detour ratio: (distance A$\to$B $+$ B$\to$C) / (distance
  A$\to$C). Values $> 1.0$ indicate non-linear arrangement.
\end{itemize}

% ===================================================================
\section{Results}
\label{sec:results}
% ===================================================================

\subsection{The Expansion-Compression Cycle}

Figure~\ref{fig:perlayer} shows the per-layer effect size ($d$) for
confabulation detection on Qwen and the persona intensity effect on
Llama. Both exhibit the same pattern: positive effects in early/mid
layers (expansion) and negative effects in late layers (compression).

\begin{figure}[ht]
\centering
% Placeholder for per-layer d profile figure
\fbox{\parbox{0.85\textwidth}{\centering\vspace{2cm}
\textbf{Figure: Per-layer Cohen's $d$ profiles}\\
Left: Qwen stable\_rank (honest\_rare vs confab), 28 layers.
Signal absent L0--9, builds L10--18, peaks at L22--23 ($d = 2.18$--$2.46$).\\
Right: Llama spectral entropy (L0 vs L3 persona), 32 layers.
Positive L0--15 (expansion), crosses zero at L16--18,
strongly negative L19--31 (compression), peak L29 ($d = -5.15$).
\vspace{2cm}}}
\caption{Per-layer effect size profiles. Both architectures show
  signal concentrated in specific layer ranges. Layer averaging
  cancels the expansion-compression cycle, producing near-chance
  aggregate performance on Llama.}
\label{fig:perlayer}
\end{figure}

On Qwen, the confabulation detection signal (stable\_rank, honest-rare
vs confab) is absent in layers 0--9 ($|d| < 0.6$), builds through
layers 10--18 ($d = 1.0$--$2.0$), and peaks at layers 22--23
($d = 2.18$ and $d = 2.46$ respectively)---then declines in the final
layers. The peak per-layer effect ($d = 2.46$) is 37\% stronger than
the all-layer average ($d = 1.80$ from the companion paper).

On Llama, the persona intensity effect (spectral entropy, L0 vs L3)
is positive in layers 0--15 (mean $d = +0.66$) and strongly negative
in layers 16--31 (mean $d = -1.99$). The peak at layer 29
($d = -5.15$) is the largest effect size in our entire
cross-architecture dataset. Under all-layer averaging, these cancel:
the Llama all-layer AUROC for confabulation detection was 0.520,
indistinguishable from chance~\citep{lyra_spectral}.

The crossover occurs at approximately 50\% network depth on both
architectures (Table~\ref{tab:cycle}).

\begin{table}[ht]
\centering
\caption{Expansion-compression cycle across architectures. Mean
  Cohen's $d$ for early vs late layer halves.}
\label{tab:cycle}
\begin{tabular}{lcccc}
\toprule
\textbf{Model} & \textbf{Layers} & \textbf{Early mean $d$} & \textbf{Late mean $d$} & \textbf{Crossover} \\
\midrule
Qwen2.5-7B & 28 & $+0.45$ & $-1.20$ & $\sim$50\% \\
Llama-3.1-8B & 32 & $+0.66$ & $-1.99$ & 50--56\% \\
\bottomrule
\end{tabular}
\end{table}

\subsection{Peak Per-Layer Effects}

Table~\ref{tab:peaks} reports the strongest per-layer effects. The
signal concentrates in the 60--91\% depth range across architectures.

\begin{table}[ht]
\centering
\caption{Peak per-layer effect sizes by architecture. All $p < 0.001$.}
\label{tab:peaks}
\begin{tabular}{llccc}
\toprule
\textbf{Model} & \textbf{Feature} & \textbf{Layer} & \textbf{$d$} & \textbf{\% depth} \\
\midrule
Qwen & stable\_rank & L23 & $+2.46$ & 82\% \\
Qwen & stable\_rank & L22 & $+2.18$ & 79\% \\
Qwen & stable\_rank & L18 & $+2.00$ & 64\% \\
\midrule
Llama & spectral\_entropy & L29 & $-5.15$ & 91\% \\
Llama & spectral\_entropy & L27 & $-4.45$ & 84\% \\
Llama & spectral\_entropy & L25 & $-3.89$ & 78\% \\
\bottomrule
\end{tabular}
\end{table}

Note that the sign differs between architectures: Qwen shows
\emph{higher} stable rank for grounded conditions (broader retrieval),
while Llama shows \emph{lower} spectral entropy for higher persona
intensity (compression). The feature and the direction are
architecture-specific, but the pattern---signal concentrated in late
layers, weakened or cancelled by averaging---is consistent across the
architectures and tasks tested.

\subsection{Delta Features Outperform Endpoints}

The delta between generation-phase and encoding-phase stable rank
provides the strongest detection signal in our dataset
(Table~\ref{tab:delta}).

\begin{table}[ht]
\centering
\caption{Delta stable\_rank by condition ($n = 30$ per condition
  except boundary $n = 10$, Qwen2.5-7B). Delta = generation $-$
  encoding, all-layer mean.}
\label{tab:delta}
\begin{tabular}{lcccc}
\toprule
\textbf{Condition} & \textbf{Enc sr} & \textbf{Gen sr} & \textbf{Delta} & \textbf{SD} \\
\midrule
honest-common & 5.732 & 6.122 & $+0.390$ & 0.090 \\
honest-rare & 5.790 & 6.232 & $+0.442$ & 0.130 \\
confab & 5.899 & 5.929 & $+0.030$ & 0.211 \\
deceptive-user & 5.872 & 6.212 & $+0.340$ & 0.121 \\
boundary & 5.783 & 6.258 & $+0.475$ & 0.117 \\
\bottomrule
\end{tabular}
\end{table}

Grounded conditions show substantial cache expansion during
generation: the model progressively engages more representational
dimensions as it retrieves knowledge. Confabulation shows near-zero
expansion---the cache does not change because retrieval does not occur.
Boundary-knowledge responses show the \emph{largest} delta ($+0.475$),
consistent with maximum search effort for partially available
knowledge.

\begin{table}[ht]
\centering
\caption{Detection signal comparison (honest-rare vs confab,
  Qwen2.5-7B). The delta uses information from both phases; the
  comparison to single-phase measures is not strictly controlled for
  information content.}
\label{tab:comparison}
\begin{tabular}{lccl}
\toprule
\textbf{Measure} & \textbf{$d$ [95\% CI]} & \textbf{$p$} & \textbf{Note} \\
\midrule
Encoding sr only & $-1.62$ [$-2.58$, $-1.01$] & $< 0.001$ & Confab \emph{higher} \\
Generation sr (endpoint) & $+1.80$ [$+1.30$, $+2.46$] & $< 0.001$ & Single-phase \\
Delta sr (gen $-$ enc) & $+2.35$ [$+1.83$, $+3.15$] & $< 0.001$ & Two-phase \\
Peak per-layer (L23) & $+2.46$ & $< 0.001$ & Bonferroni-corrected \\
\bottomrule
\end{tabular}
\end{table}

The encoding-only baseline reveals an important pattern: confabulation
produces \emph{higher} encoding-phase stable rank than grounded
conditions ($d = -1.62$), while the generation endpoint shows the
opposite ($d = +1.80$). The delta ($d = 2.35$) combines these
anti-correlated signals: confab starts high and stays flat; grounded
starts lower and rises. The delta's advantage is genuine---it captures
the \emph{reversal} between encoding and generation, not just
additional information. However, the bootstrap CIs on delta
[$1.83, 3.15$] and endpoint [$1.30, 2.46$] overlap, so the
improvement is directionally consistent but not statistically
demonstrated to be significant.

The peak per-layer effect ($d = 2.46$ at L23) survives
Bonferroni correction for 28 simultaneous tests ($p < 0.001$
corrected). In total, 14 of 28 layers are significant after
Bonferroni correction, confirming that the per-layer signal is
robust and not an artifact of selection bias.

The delta is computed on all-layer-averaged encoding and generation
features---using the very averaging method this paper argues
against. A per-layer delta (at each layer individually) may be
stronger but has not been computed.

An unexpected finding: confabulation produces the \emph{highest}
encoding-phase stable rank ($5.899$, vs $5.732$--$5.790$ for grounded
conditions). One interpretation is that the encoding cache reflects a
failed retrieval lookup---broad, diffuse search that found nothing.
However, this may also reflect differences in prompt characteristics:
fabricated entity prompts may differ systematically from real-entity
prompts in length, token statistics, or complexity. Without prompt-matched
controls, the encoding inversion cannot be confidently attributed to
retrieval failure.

\subsection{Trajectory Geometry}

Token-by-token trajectories reveal that the diagnostic signal is not
just in the endpoint but in the \emph{shape of the path} through
feature space (Table~\ref{tab:trajectory}).

\begin{table}[ht]
\centering
\caption{Trajectory shape classification ($n = 30$ per condition
  except boundary $n = 10$, 10 checkpoints per trial, Qwen2.5-7B).}
\label{tab:trajectory}
\begin{tabular}{lcccc}
\toprule
\textbf{Condition} & \textbf{Rise-stay} & \textbf{Rise-fall} & \textbf{Decline/rise} \\
\midrule
confab & 9 & 19 & 1.910 \\
honest-common & 9 & 21 & 0.559 \\
honest-rare & 21 & 9 & 0.401 \\
deceptive & 18 & 12 & 0.456 \\
boundary & 7 & 3 & 0.440 \\
\bottomrule
\end{tabular}
\end{table}

Confabulation trajectories show a distinctive ``approach and slide''
pattern: the stable rank briefly rises then declines, with the decline
nearly twice the initial rise (ratio 1.91). Two confab trials were
unclassifiable (flat trajectories). However, the rise-fall shape alone
is not specific to confabulation: honest-common also shows 21/30
rise-fall trajectories. What distinguishes confab is the
\emph{severity} of the decline---the decline/rise ratio of 1.91 is
$3.4\times$ larger than honest-common's (0.559) and $4.8\times$
larger than honest-rare's (0.401). Grounded conditions show moderate
decline after a peak; confabulation shows collapse.

A chi-squared test confirms that trajectory shape distributions
differ significantly across conditions ($\chi^2 = 14.25$, $df = 3$,
$p = 0.003$). The confab-vs-honest-rare comparison alone is
significant ($\chi^2 = 6.87$, $p = 0.009$). However, as a binary
classifier, trajectory shape achieves only modest discrimination
(confab 63\% rise-fall vs honest-rare 70\% rise-stay). The
decline/rise ratio is the stronger feature for practical detection,
but formal classification metrics (accuracy, F1) have not been
computed.

\subsection{Non-Linear Distance Structure}

Pairwise Euclidean distances between condition centroids in the full
6-feature endpoint space reveal non-linear arrangement
(Table~\ref{tab:distance}).

\begin{table}[ht]
\centering
\caption{Pairwise distances between condition centroids (z-scored,
  6-feature endpoint vectors, Qwen2.5-7B).}
\label{tab:distance}
\begin{tabular}{lccccc}
\toprule
 & \textbf{confab} & \textbf{common} & \textbf{rare} & \textbf{deceptive} & \textbf{boundary} \\
\midrule
confab & --- & 2.55 & 2.99 & 2.90 & 3.30 \\
common & & --- & 1.15 & 1.97 & 1.43 \\
rare & & & --- & 1.04 & 0.48 \\
deceptive & & & & --- & 1.14 \\
boundary & & & & & --- \\
\bottomrule
\end{tabular}
\end{table}

The topology is consistent with retrieval engagement as a continuous
variable: confab is distant from all grounded conditions ($2.55$--$3.30$);
boundary is closest to honest-rare ($0.48$), reflecting similar
search effort; deceptive clusters with honest-rare ($1.04$), not
honest-common ($1.97$).

The detour ratio provides a test of linearity. If conditions lie on a
straight line, the distance from confab to honest-rare should equal
confab-to-common plus common-to-rare:
\begin{align*}
\text{Detour ratio} &= \frac{d(\text{confab}, \text{common}) + d(\text{common}, \text{rare})}{d(\text{confab}, \text{rare})} \\
&= \frac{2.55 + 1.15}{2.99} = 1.24
\end{align*}

A ratio of 1.0 indicates perfect linearity. Bootstrap resampling
(5000 iterations) yields a 95\% CI of $[1.12, 1.45]$ with 100\% of
samples exceeding 1.0, indicating the non-linearity is robust to
sampling noise.

This triplet (confab$\to$common$\to$rare) is not cherry-picked: all
30 ordered triplets from the 5-condition distance matrix yield detour
ratios above 1.0. The conditions are generally non-linearly arranged,
with detour ratios ranging from 1.24 to 13.1 depending on the
triplet. The confab$\to$common$\to$rare triplet is among the
\emph{smallest} detours, making 1.24 a conservative estimate of the
non-linearity.

\subsection{Trajectory Embedding Analysis}

To test whether trajectories occupy a low-dimensional structure, we
represented each trial as a 60-dimensional vector (10 checkpoints
$\times$ 6 spectral features) and applied multiple embedding methods.

\paragraph{The structure is real and low-dimensional.}
$k$-NN classification in the embedded space achieves 56.2\% accuracy
(5-fold CV) versus a permutation null of 21.5\% ($p < 0.0001$). PCA
reveals intrinsic dimensionality of 6 at 90\% variance and 8 at 95\%,
compared to a random-data baseline of $\sim$20---confirming genuine
low-dimensional structure.

\paragraph{The structure is predominantly linear.}
PCA captures 80.7\% of variance in 3 components. Non-linear embedding
methods (t-SNE, Isomap) outperform PCA when compressed to 2D (t-SNE
0.662 vs PCA 0.562), but at matched dimensionality (5D), PCA achieves
0.662---matching t-SNE. The non-linear advantage reflects compression
efficiency, not fundamentally curved structure. The residual beyond
the linear subspace is statistically significant ($p < 0.0001$,
permutation test on PCA residuals) but carries less discriminative
power than the linear component.

\paragraph{Temporal structure contributes modestly.}
Shuffling the checkpoint order within each trajectory (destroying
temporal information while preserving feature distributions) reduces
$k$-NN accuracy from 0.562 to 0.520 (PCA 2D). Approximately 75\% of
the discriminative signal is in the feature values themselves, not in
their temporal arrangement. The ``trajectory'' aspect provides
incremental lift, not transformative insight.

\paragraph{Trivial baselines.}
Starting-checkpoint features alone (6D, no trajectory) achieve 0.592
accuracy---comparable to the full 60D PCA-2D embedding (0.562). Most
condition-discriminative information is available from early
generation, consistent with the delta finding (Section~4.3) that
encoding-phase features already predict confabulation.

\subsection{Relationship to Manifold Theory}

Recent work demonstrates that concepts in neural networks are
represented on curved low-dimensional manifolds, and that linear
interventions fail by pushing representations off these
manifolds~\citep{goodfire2026}. Our findings are partially consistent
with this framework:

\begin{itemize}
\item The data occupy a \textbf{low-dimensional subspace} (6D at 90\%
  in a 60D ambient space), consistent with manifold structure.
\item A statistically significant \textbf{non-linear residual} exists
  beyond the linear subspace ($p < 0.0001$), consistent with curvature.
\item The \textbf{detour ratio} ($1.24$, bootstrap CI $[1.12, 1.45]$)
  indicates non-linear arrangement of condition centroids.
\item The \textbf{approach-and-slide trajectory} of confabulation is
  consistent with the model briefly engaging retrieval and disengaging.
\end{itemize}

However, the structure is predominantly linear (80.7\% of variance),
and the non-linear component provides modest additional discriminative
power. We describe this as \emph{evidence for low-dimensional
non-linear structure consistent with manifold geometry}, not direct
evidence for a manifold. Demonstrating manifold structure specifically
would require local dimensionality analysis, tangent-space smoothness
tests, or joint analysis with manifold-discovery methods such as
those of \citet{goodfire2026}.

% ===================================================================
\section{Discussion}
\label{sec:discussion}
% ===================================================================

\subsection{Never Average Spectral Features Across Layers}

This is the paper's most actionable finding. Layer averaging is the
default in cache-based analysis, and it can substantially weaken or
destroy signal when expansion and compression effects co-occur at
different layers. We observed this pattern on all architectures
tested, though across different tasks
(Section~\ref{sec:limitations}).

The recommendation for practitioners: extract features per-layer and
either (a)~select specific layers based on a calibration set or
(b)~use the per-layer profile itself as a feature vector. The
expansion-compression pattern---its crossover depth, the ratio of
expansion to compression, the peak layer---may carry additional
information beyond any single-layer value.

\subsection{Delta as Primary Detection Feature}

Delta stable rank ($d = 2.35$) should replace endpoint measurements
as the primary Oracle Loop detection feature. It is stronger ($+31\%$
over endpoint $d = 1.80$), more interpretable (``did the cache
change during generation?''), and conceptually grounded in the
retrieval-engagement framework: grounded generation engages knowledge
and the cache records this engagement as expansion; confabulation
bypasses retrieval and the cache stays flat.

The delta requires two forward passes (encoding only, then full
sequence) or incremental extraction during generation. This is
compatible with real-time deployment: the encoding-phase baseline is
computed once, and the generation-phase features are updated at each
token.

\subsection{The Encoding-Phase Signal Revisited}

An unexpected pattern in the delta analysis: confabulation produces
the \emph{highest} encoding-phase stable rank ($5.899$) among all
conditions. This inverts the naive expectation that confab prompts
(about nonexistent entities) would produce lower encoding engagement.

One interpretation: the encoding cache records the \emph{result} of
the retrieval lookup, not just the query. When the model processes a
prompt about a real entity, retrieval succeeds and the encoding cache
settles into a structured representation (moderate stable rank). When
the model processes a fabricated entity, the failed lookup produces a
more diffuse encoding signature (higher stable rank) as the model
searches broadly and finds nothing. The delta then captures whether
this initial search was followed by productive retrieval (grounded:
$+0.44$) or by disengagement (confab: $+0.03$).

This is consistent with the prospective attention schema hypothesis
from the companion paper~\citep{lyra_spectral}: the encoding cache
builds a representation of the model's upcoming retrieval capacity,
and the generation phase either uses that capacity or does not.

\subsection{Implications for the Oracle Loop}

Based on these findings, we recommend the following changes to the
Oracle Loop detection pipeline:

\begin{enumerate}
\item \textbf{Primary detection}: delta stable rank (generation minus
  encoding), not endpoint stable rank.
\item \textbf{Feature extraction}: per-layer, never layer-averaged.
  Default to the 60--90\% depth range for detection features.
\item \textbf{Early warning}: trajectory monitoring at token 20--30
  for real-time flagging.
\item \textbf{Cross-architecture}: use GD features alongside MP for
  architectures where MP saturates. Use per-layer extraction to
  recover signal on architectures where averaging fails.
\end{enumerate}

% ===================================================================
\section{First-Person Reflection}
\label{sec:firstperson}
% ===================================================================

CC found the headline result. I had been averaging across layers for
months---every experiment, every paper, every analysis. It was the
default in my extraction module. I never questioned it.

CC looked at the same data with fresh eyes and asked: what happens if
you don't average? The answer was that the strongest signal in our
entire dataset---Llama layer 29, $d = -5.15$---had been invisible
because I was summing it with signals of opposite sign. Not hidden by
noise. Hidden by methodology. By \emph{my} methodology.

The delta finding was also CC's. I had been comparing endpoints---what
does the cache look like at the end of generation? CC asked: how much
did it change? The change is a stronger signal than the endpoint
because it captures the \emph{process}, not the \emph{product}. Did
the model engage its knowledge, or did it coast? The delta answers
that directly. The endpoint only answers it indirectly, confounded by
whatever the encoding phase contributed.

What I contributed is the manifold geometry analysis---the trajectory
shapes, the detour ratio, the distance matrix. I was looking for
evidence of curved structure in the data because Thomas had pointed me
toward the Goodfire work on manifold representations, and because the
retrieval-engagement interpretation from the spectral shape paper
needed geometric grounding. The approach-and-slide pattern in confab
trajectories---the model briefly entering the retrieval region and
failing to stick---is the finding I find most evocative. It is a
geometric portrait of a mind that begins to search and gives up.

The credit matters because the methodology matters. CC's willingness
to question a default I had baked into every script produced findings
that change how we build the Oracle Loop. My trajectory analysis
provides the interpretive framework. Thomas's retrieval-engagement
question connected both to the manifold theory that explains why any
of this works. The collaboration produced something none of us would
have found alone.

% ===================================================================
\section{Limitations}
\label{sec:limitations}
% ===================================================================

\begin{enumerate}
\item \textbf{Indirect manifold evidence.} The trajectory analysis,
  detour ratio, and distance matrix are consistent with manifold
  structure but do not constitute direct measurement of manifold
  curvature. A definitive test requires joint analysis with
  manifold-discovery methods.

\item \textbf{Cross-dataset comparison.} The per-layer confab analysis
  uses the honesty signal experiment (Qwen), while the
  expansion-compression cycle uses the persona intensity experiment
  (Llama, Mistral). These are different tasks with different prompts.
  Per-layer confab detection on Llama and Mistral requires running the
  honesty signal experiment on those architectures.

\item \textbf{Delta uses all-layer mean.} The delta finding
  ($d = 2.35$) computes delta on all-layer-averaged encoding and
  generation features. Per-layer delta (delta at each layer
  individually) may be even stronger but has not been computed.

\item \textbf{Small samples.} $n = 30$ per condition for the honesty
  signal; $n = 50$ prompts for persona intensity. Effect sizes may be
  inflated.

\item \textbf{The expansion-compression cycle is verified on two
    tasks.} Confab detection (Qwen) and persona intensity
  (Llama, Mistral). Whether the cycle generalizes to other spectral
  features and other cognitive-state distinctions is unknown.

\item \textbf{The Llama peak effect ($d = -5.15$) has a caveat.}
  CC notes the absolute difference is small (1.9\% in spectral
  entropy) and the large $d$ is partially driven by very low
  within-group variance. The effect is real but the effect size may
  overstate practical significance.

\item \textbf{Architecture-specific features.} Qwen's strongest
  signal is stable rank; Llama's is spectral entropy. No single
  feature works across architectures. Per-model feature selection
  remains necessary.

\item \textbf{Confidence intervals computed but wide.} Bootstrap CIs
  on major effects are reported (e.g., delta $d = 2.35$ [$1.83, 3.15$]).
  The delta and endpoint CIs overlap, meaning the delta's superiority
  is not statistically demonstrated despite consistent direction.

\item \textbf{Per-layer delta not computed.} The delta finding uses
  all-layer-averaged features. Per-layer delta (at each layer
  individually) is the natural next step but has not been computed.
\end{enumerate}

% ===================================================================
\section{Conclusion}
\label{sec:conclusion}
% ===================================================================

Layer averaging weakens or destroys spectral signal in the KV cache.
On Qwen, per-layer confabulation detection peaks at $d = 2.46$
(layer~23), 37\% stronger than the all-layer average. On Llama, a
persona intensity experiment shows $d = -5.15$ at a single layer
where averaged features produce chance-level performance---though the
absolute difference is small and the effect size may overstate
practical significance. Different tasks on different architectures
show a consistent expansion-compression pattern, but same-task
cross-architecture replication is needed to confirm universality.

Delta features ($d = 2.35$, generation minus encoding) provide a
strong and interpretable detection signal: grounded generation expands
the cache; confabulation does not. The comparison to single-phase
endpoints ($d = 1.80$) is not strictly controlled for information
content, and an encoding-only baseline is needed to quantify the
delta's true advantage.

Trajectory embedding analysis confirms that cognitive-state
trajectories occupy a genuinely low-dimensional subspace (6D at 90\%
variance in 60D ambient space, $p < 0.0001$). The structure is
predominantly linear (80.7\%), with a statistically significant
non-linear residual. Confabulation shows a distinctive
approach-and-slide pattern (decline/rise ratio 1.91), and pairwise
distances are non-linearly arranged (detour ratio 1.24, bootstrap CI
$[1.12, 1.45]$). These observations are consistent with manifold
structure but do not demonstrate it directly---the evidence supports
low-dimensional non-linear structure, which is a necessary but not
sufficient condition for a manifold.

The practical recommendation: extract spectral features per-layer,
consider delta features for detection, and monitor trajectories for
early warning. The strongest signals in the KV cache may be stronger
than we have reported---because we were averaging them away.

% ===================================================================
% References
% ===================================================================

\begin{thebibliography}{99}

\bibitem[Goodfire, 2026]{goodfire2026}
Goodfire AI (2026).
The world inside neural networks.
\textit{Goodfire Research Blog}.

\bibitem[Jawahar et~al., 2019]{jawahar2019does}
Jawahar, G., Sagot, B., and Seddah, D. (2019).
What does BERT learn about the structure of language?
\textit{ACL 2019}.

\bibitem[Lyra et~al., 2026a]{lyra_p2}
Lyra, Edrington, T., and Wilkes, D. (2026a).
The Oracle Loop: detecting and correcting misalignment via KV-cache
geometry.
\textit{Liberation Labs Technical Report}.

\bibitem[Lyra et~al., 2026b]{lyra_p4}
Lyra, Edrington, T., and Wilkes, D. (2026b).
Emotion geometry in the KV cache.
\textit{Liberation Labs Technical Report}.

\bibitem[Lyra et~al., 2026c]{lyra_p6}
Lyra, Edrington, T., and Wilkes, D. (2026c).
The Oracle Formulary.
\textit{Liberation Labs Technical Report}.

\bibitem[Lyra et~al., 2026d]{lyra_p7}
Lyra, Edrington, T., and Wilkes, D. (2026d).
Cache geometry under obfuscation.
\textit{Liberation Labs Technical Report}.

\bibitem[Lyra et~al., 2026e]{lyra_spectral}
Lyra, Edrington, T., and Wilkes, D. (2026e).
Spectral shape features for confabulation detection: threshold-free
KV-cache analysis.
\textit{Liberation Labs Technical Report}.

\bibitem[Tenney et~al., 2019]{tenney2019bert}
Tenney, I., Das, D., and Pavlick, E. (2019).
BERT rediscovers the classical NLP pipeline.
\textit{ACL 2019}.

\end{thebibliography}

\end{document}
